\documentclass[conference]{IEEEtran}
\IEEEoverridecommandlockouts

% ============================================================
% PACKAGES
% ============================================================
\usepackage{cite}
\usepackage{amsmath,amssymb,amsfonts}
\usepackage{algorithmic}
\usepackage{graphicx}
\usepackage{textcomp}
\usepackage{xcolor}
\usepackage{booktabs}
\usepackage{multirow}
\usepackage{array}
\usepackage{hyperref}
\usepackage{float}
\usepackage{enumitem}
\usepackage{subcaption}
\usepackage{balance}

\def\BibTeX{{\rm B\kern-.05em{\sc i\kern-.025em b}\kern-.08em
    T\kern-.1667em\lower.7ex\hbox{E}\kern-.125emX}}

\begin{document}

% ============================================================
% TITLE
% ============================================================
\title{Graph-Enhanced Ensemble Learning for Healthcare Provider Fraud Detection: A Neo4j-Augmented Approach with Overfitting-Aware Regularization}

\author{
\IEEEauthorblockN{Ashish Kumar}
\IEEEauthorblockA{
\textit{Department of Data Science} \\
\textit{Capstone Project} \\
Email: ashish@example.com}
}

\maketitle

% ============================================================
% ABSTRACT
% ============================================================
\begin{abstract}
Healthcare fraud costs the United States an estimated \$100--\$300 billion annually, making automated detection a critical priority for the Centers for Medicare \& Medicaid Services (CMS). This paper presents a production-grade, graph-enhanced ensemble learning pipeline for provider-level fraud detection that combines tabular claim aggregations with Neo4j-derived network features. We engineer 58 domain-specific tabular features and 12 graph-based features capturing provider network topology, shared-beneficiary overlap, diagnosis concentration, and beneficiary exclusivity patterns. Through an iterative model improvement process, we demonstrate how explicit overfitting diagnosis---comparing train versus validation scores across folds---guides regularization tuning. Our improved calibrated stacking ensemble of XGBoost, LightGBM, and CatBoost achieves a PR-AUC of 0.7042 and recall of 76.5\% on out-of-fold predictions across 5,410 Medicare providers, while reducing the train--validation PR-AUC gap by up to 72\% compared to the initial configuration. We detail the complete progression from an initial model exhibiting significant overfitting (train PR-AUC = 1.0, validation PR-AUC = 0.70) through aggressive regularization, feature selection, repeated stratified cross-validation, adversarial validation, and calibration assessment that yields a robust, defensible system. The pipeline is designed for deployment on Windows with NVIDIA GPU acceleration and Neo4j Community Edition.
\end{abstract}

\begin{IEEEkeywords}
healthcare fraud detection, ensemble learning, graph features, Neo4j, overfitting diagnosis, XGBoost, LightGBM, CatBoost, class imbalance, Medicare claims
\end{IEEEkeywords}

% ============================================================
% I. INTRODUCTION
% ============================================================
\section{Introduction}

Healthcare fraud remains one of the most financially damaging forms of white-collar crime in the United States. The National Health Care Anti-Fraud Association (NHCAA) estimates that fraud accounts for 3--10\% of total healthcare expenditures, translating to \$100--\$300 billion in annual losses \cite{nhcaa2020}. The Centers for Medicare \& Medicaid Services (CMS) processes over 1 billion claims annually, making manual audit infeasible and motivating the development of machine learning systems for automated fraud detection \cite{bauder2017}.

Traditional rule-based systems, while interpretable, suffer from high false-positive rates and inability to detect novel fraud schemes \cite{herland2018}. Machine learning approaches have demonstrated superior detection performance, particularly gradient-boosted tree ensembles on tabular claims data \cite{johnson2023}. However, most published work treats providers as independent entities, ignoring the rich network structure inherent in healthcare claims---providers share patients, employ overlapping physician networks, and exhibit billing patterns that are best understood in relational context.

A critical but often under-reported challenge in healthcare fraud detection is \textbf{overfitting}. With small provider-level datasets (thousands, not millions, of rows), tree-based models can easily memorize training data, producing inflated performance estimates that do not generalize. Many published systems report only validation metrics without comparing them against training scores, making it impossible to assess whether reported performance is trustworthy.

This paper makes the following contributions:

\begin{enumerate}[leftmargin=*]
    \item \textbf{Graph-augmented feature engineering}: We construct a provider--beneficiary--physician--diagnosis knowledge graph in Neo4j and extract 12 network features capturing degree centrality, shared-entity overlap, concentration indices, and beneficiary exclusivity.
    \item \textbf{Explicit overfitting diagnosis and mitigation}: We introduce train-vs-validation gap tracking across all CV folds, demonstrating how an initial model with perfect train scores (ROC-AUC = 1.0) is systematically improved through aggressive regularization, reducing the PR-AUC gap from 0.295 to 0.054.
    \item \textbf{Improved stacking ensemble}: A calibrated ensemble with stronger regularization achieves higher recall (76.5\% vs 72.9\%) and PR-AUC (0.7042 vs 0.7010) while being substantially less overfit.
    \item \textbf{Comprehensive diagnostics}: We include adversarial validation, calibration assessment (Brier score, ECE), repeated stratified CV, and importance-based feature pruning as standard pipeline components.
    \item \textbf{Production-ready implementation}: Complete, modular Python codebase with Neo4j integration, GPU support, and Windows compatibility.
\end{enumerate}

The remainder of this paper is organized as follows: Section~\ref{sec:related} reviews related work. Section~\ref{sec:methodology} describes the dataset, feature engineering, graph construction, and modeling pipeline. Section~\ref{sec:results} presents experimental results including the overfitting analysis. Section~\ref{sec:discussion} discusses findings and limitations. Section~\ref{sec:conclusion} concludes.

% ============================================================
% II. RELATED WORK
% ============================================================
\section{Related Work}
\label{sec:related}

\subsection{Machine Learning for Healthcare Fraud Detection}

Bauder and Khoshgoftaar \cite{bauder2017} pioneered the use of the CMS Medicare Provider Utilization and Payment Public Use File (PUF) combined with the Office of Inspector General's List of Excluded Individuals and Entities (LEIE) for supervised fraud detection. They evaluated multiple classifiers including Random Forest, Logistic Regression, and na\"ive Bayes, finding that ensemble tree methods consistently outperformed linear models on the highly imbalanced provider fraud task.

Herland et al. \cite{herland2018} extended this line of research by systematically combining CMS Part B, Part D, and DMEPOS public use files with LEIE labels across multiple years (2012--2016). Their study in the \textit{Journal of Big Data} demonstrated that multi-source data fusion significantly improved fraud detection performance, reporting ROC-AUC values ranging from 0.92 to 0.97 depending on the data combination and year.

Johnson and Khoshgoftaar \cite{johnson2023} provided a comprehensive survey and empirical comparison of Medicare fraud detection approaches in \textit{SN Computer Science}, identifying gradient-boosted decision trees (XGBoost, LightGBM, CatBoost) as the most reliable family for tabular Medicare claims data. They highlighted persistent challenges: extreme class imbalance (fraud rates of 1--10\%), label noise, temporal distribution shift, and the risk of data leakage in provider-level aggregation pipelines.

\subsection{Graph-Based Fraud Detection}

Network analysis has been extensively applied in financial fraud detection \cite{pourhabibi2020}, but its application to healthcare provider fraud remains relatively underexplored. Liu et al. \cite{liu2021} demonstrated that provider--patient bipartite graphs could reveal suspicious referral patterns invisible to tabular features alone. Van Vlasselaer et al. \cite{vanvlasselaer2017} showed that propagation-based features in transaction networks improved credit card fraud detection, a methodology transferable to healthcare claim networks.

Graph database technologies, particularly Neo4j, have been adopted in fraud detection pipelines for their ability to efficiently traverse relationship-heavy data \cite{neo4j2023}. Our work bridges the gap between graph-based visualization and quantitative feature extraction by systematically deriving features from a Neo4j knowledge graph and demonstrating their additive value in a gradient-boosted ensemble.

\subsection{Ensemble Methods and Calibration}

Stacking ensembles that combine diverse base learners through a meta-learner have shown consistent improvements over single models in tabular data settings \cite{wolpert1992}. For fraud detection, the combination of XGBoost, LightGBM, and CatBoost provides complementary inductive biases \cite{chen2016, ke2017, prokhorenkova2018}. Probability calibration is particularly important in fraud detection, where downstream decision systems require well-calibrated risk scores for cost-sensitive thresholding \cite{niculescu2005}.

% ============================================================
% III. METHODOLOGY
% ============================================================
\section{Methodology}
\label{sec:methodology}

\subsection{Dataset Description}

We use the CMS Healthcare Provider Fraud Detection dataset, which contains Medicare claims data with provider-level fraud labels. The dataset comprises eight CSV files:

\begin{itemize}[leftmargin=*]
    \item \textbf{Provider labels}: 5,410 training providers (506 fraudulent, 9.35\% positive rate) and 1,353 test providers.
    \item \textbf{Beneficiary data}: 138,557 training and 63,969 test beneficiary records with demographics and chronic condition indicators.
    \item \textbf{Inpatient claims}: 40,475 training and 9,552 test hospitalization records with diagnosis codes (up to 10), procedure codes (up to 6), reimbursement amounts, and physician identifiers.
    \item \textbf{Outpatient claims}: 517,738 training and 125,842 test outpatient visit records.
\end{itemize}

After merging inpatient and outpatient claims with beneficiary data, the training set contains 558,211 claim-level rows aggregated to the provider level for modeling.

\subsection{Data Preprocessing}

\subsubsection{Claim Merging and Type Tagging}
Inpatient and outpatient claims are separately merged with beneficiary records on \texttt{BeneID}, then concatenated with a binary \texttt{IsInpatient} indicator.

\subsubsection{Date Parsing and Duration Computation}
All date columns are parsed with \texttt{coerce} error handling. Claim duration is computed as $(\texttt{DischargeDt} - \texttt{AdmissionDt}).days$, floored at zero.

\subsubsection{Chronic Condition Encoding}
Eleven chronic condition indicators are binary-encoded by mapping CMS value ``1'' to 1 and ``2'' to 0. A composite \texttt{ChronicCount} feature sums across all conditions per beneficiary.

\subsubsection{Age and Mortality Features}
Beneficiary age is computed relative to a fixed reference date. A binary \texttt{IsDead} indicator flags deceased beneficiaries.

\subsection{Provider-Level Feature Engineering}

All features are aggregated from claim-level records to the provider level. We engineer 58 tabular features organized into nine categories, summarized in Table~\ref{tab:features}.

\begin{table}[t]
\caption{Provider-Level Tabular Feature Categories}
\label{tab:features}
\centering
\small
\begin{tabular}{@{}lcp{4.2cm}@{}}
\toprule
\textbf{Category} & \textbf{Count} & \textbf{Key Features} \\
\midrule
Claim Volume & 5 & TotalClaims, InpatientRatio, ClaimsPerBeneficiary \\
Financial & 8 & TotalReimbursed, MeanReimbursed, StdReimbursed, ReimbursedCV \\
Beneficiary & 6 & UniqueBeneficiaries, MeanAge, FracDead, MeanChronicCount \\
Chronic Rates & 12 & Per-condition prevalence rates \\
Dx/Px Diversity & 8 & UniqueDiagCodes, DiagDiversityRatio, ProcDiversityRatio \\
Physician & 5 & TotalUniquePhysicians, UniqueAttending, UniqueOperating \\
Temporal & 3 & MeanClaimDuration, MaxClaimDuration, ClaimDateSpanDays \\
Geographic & 4 & UniqueStates, UniqueCounties, MaleRatio, UniqueRaces \\
Anomaly Ratios & 3 & ZeroDeductibleFrac, MaxSameDayClaimsPerBene \\
\bottomrule
\end{tabular}
\vspace{-2mm}
\end{table}

\subsection{Neo4j Graph Construction and Feature Extraction}

\subsubsection{Graph Schema}
We construct a heterogeneous knowledge graph in Neo4j with five node types: \texttt{Provider}, \texttt{Beneficiary}, \texttt{Physician}, \texttt{DiagnosisCode}, and \texttt{ProcedureCode}; and six relationship types: \texttt{TREATED}, \texttt{EMPLOYED}, \texttt{VISITED}, \texttt{WORKS\_AT}, \texttt{USED\_DIAGNOSIS}, and \texttt{USED\_PROCEDURE}. Nodes and relationships are ingested via batched parameterized Cypher queries.

\subsubsection{Graph-Derived Features}
We extract 12 features from the graph:

\textbf{Degree features} (4): Distinct counts of beneficiaries, physicians, diagnosis codes, and procedure codes per provider.

\textbf{Shared-entity overlap features} (5): For each provider pair sharing at least one beneficiary:
\begin{equation}
    \text{shared\_bene\_partners}_i = |\{j : |B_i \cap B_j| > 0, j \neq i\}|
\end{equation}
Similarly, shared physician overlap counts are computed. These features capture network embeddedness---fraudulent providers may operate in clusters sharing patient pools.

\textbf{Concentration features} (2): Herfindahl-Hirschman Index (HHI) for diagnosis and procedure code distributions:
\begin{equation}
    \text{HHI}_{\text{diag},i} = \sum_{d \in D_i} \left(\frac{n_{i,d}}{N_i}\right)^2
\end{equation}
High concentration (HHI $\to$ 1) indicates a provider billing a narrow set of codes, a known fraud indicator.

\textbf{Beneficiary exclusivity} (1): Fraction of a provider's beneficiaries who visit only that provider.

When Neo4j is unavailable, equivalent features are computed using offline pandas operations with set-intersection joins.

\subsection{Feature Selection Pipeline}

We apply a three-stage feature selection process:

\begin{enumerate}[leftmargin=*]
    \item \textbf{Correlation pruning} ($|r| > 0.90$): Removes one feature from each highly correlated pair. The improved threshold of 0.90 (reduced from the initial 0.95) drops 20 features, keeping 50 from the original 70.
    \item \textbf{Importance-based pruning}: After an initial XGBoost fit on one CV fold, features with importance $< 0.005$ are removed, yielding the final 49 features.
    \item \textbf{Adversarial validation}: A classifier trained to distinguish train from test rows achieves ROC-AUC = 0.488, confirming no distribution shift between splits.
\end{enumerate}

\subsection{Modeling Pipeline}

\subsubsection{Validation Strategy}
We use \textbf{3$\times$5-fold Repeated Stratified K-Fold} cross-validation (15 total evaluations per model). The repeated design provides more stable metric estimates than single-pass CV, particularly important given the small positive class (approximately 101 fraud cases per fold). Out-of-fold predictions are averaged across repeats.

\subsubsection{Class Imbalance Handling}
We use \texttt{scale\_pos\_weight} = $N_{\text{neg}} / N_{\text{pos}}$ = 9.69 for XGBoost/LightGBM and \texttt{auto\_class\_weights=Balanced} for CatBoost. We deliberately avoid SMOTE because provider-level features represent complex aggregate distributions; synthetic interpolation creates unrealistic feature combinations \cite{vandenheuvel2022}.

\subsubsection{Base Learners --- Improved Regularization}

A key contribution of this work is the \textbf{regularization-guided improvement} from an initial configuration exhibiting severe overfitting to a production-ready configuration. Table~\ref{tab:params} summarizes the parameter changes.

\begin{table}[t]
\caption{Hyperparameter Changes: Initial $\to$ Improved}
\label{tab:params}
\centering
\small
\begin{tabular}{@{}lcc@{}}
\toprule
\textbf{Parameter} & \textbf{Initial} & \textbf{Improved} \\
\midrule
\texttt{max\_depth} & 6 & \textbf{4} \\
\texttt{learning\_rate} & 0.05 & \textbf{0.03} \\
\texttt{min\_child\_weight} & 10 & \textbf{30} \\
\texttt{subsample} & 0.8 & \textbf{0.7} \\
\texttt{colsample\_bytree} & 0.7 & \textbf{0.6} \\
\texttt{reg\_alpha} (L1) & 0.5 & \textbf{2.0} \\
\texttt{reg\_lambda} (L2) & 2.0 & \textbf{5.0} \\
\texttt{gamma} & 0.3 & \textbf{1.0} \\
\texttt{num\_leaves} (LGB) & 31 & \textbf{15} \\
\texttt{l2\_leaf\_reg} (CB) & 3.0 & \textbf{10.0} \\
\bottomrule
\end{tabular}
\vspace{-2mm}
\end{table}

\textbf{XGBoost}: Histogram-based tree method with \texttt{max\_depth}=4, \texttt{learning\_rate}=0.03, \texttt{min\_child\_weight}=30, \texttt{subsample}=0.7, \texttt{colsample\_bytree}=0.6, \texttt{reg\_alpha}=2.0, \texttt{reg\_lambda}=5.0, \texttt{gamma}=1.0, \texttt{max\_delta\_step}=1. Early stopping on PR-AUC with 50-round patience.

\textbf{LightGBM}: Leaf-wise growth with \texttt{num\_leaves}=15, \texttt{max\_depth}=4, \texttt{min\_child\_samples}=40, \texttt{path\_smooth}=1.0, \texttt{reg\_alpha}=2.0, \texttt{reg\_lambda}=5.0.

\textbf{CatBoost}: Ordered boosting with \texttt{depth}=4, \texttt{l2\_leaf\_reg}=10.0, \texttt{min\_data\_in\_leaf}=40, \texttt{random\_strength}=2.0.

\subsubsection{Stacking Architecture}
Out-of-fold (OOF) predictions from the three base learners serve as meta-features for a logistic regression meta-learner (\texttt{C}=1.0, L2 penalty), wrapped in \texttt{CalibratedClassifierCV} with isotonic regression (5-fold internal CV).

\subsubsection{Threshold Optimization}
We optimize the decision threshold on OOF predictions by maximizing F1 score and additionally report precision/recall at fixed recall targets of 70\%, 80\%, and 90\%.

\subsection{Anti-Overfitting Safeguards}

Table~\ref{tab:safeguards} summarizes the ten explicit anti-overfitting mechanisms in the improved pipeline.

\begin{table}[t]
\caption{Anti-Overfitting Safeguards}
\label{tab:safeguards}
\centering
\small
\begin{tabular}{@{}p{2.8cm}p{4.8cm}@{}}
\toprule
\textbf{Safeguard} & \textbf{Threat Addressed} \\
\midrule
3$\times$5 Repeated CV & Unstable validation from small positive class \\
Early stopping (50 rounds) & Training beyond optimal generalization \\
Aggressive L1/L2 & Overly complex split decisions \\
Subsampling (0.7 row, 0.6 col) & Memorization of training examples \\
Shallow trees (depth 4) & Excessive model capacity for 5K rows \\
Correlation pruning ($r > 0.90$) & Multicollinearity inflation \\
Importance pruning & Noise features increasing variance \\
scale\_pos\_weight (no SMOTE) & Synthetic artifacts from oversampling \\
Threshold tuning on OOF & Over-optimistic in-sample thresholds \\
Adversarial validation & Undetected train/test distribution shift \\
\bottomrule
\end{tabular}
\vspace{-2mm}
\end{table}

% ============================================================
% IV. EXPERIMENTAL RESULTS
% ============================================================
\section{Experimental Results}
\label{sec:results}

\subsection{Experimental Setup}

All experiments were conducted using Python 3.10+, XGBoost 2.0+, LightGBM 4.0+, CatBoost 1.2+, and scikit-learn 1.3+. The complete improved pipeline (data loading, feature engineering, graph features, adversarial validation, feature selection, 3$\times$5-fold CV for three base learners, meta-learner training, and calibration) executes in approximately 18.5 seconds on CPU.

\subsection{Overfitting Diagnosis: Initial Model}

Our initial model used standard regularization parameters (\texttt{max\_depth}=6, \texttt{reg\_alpha}=0.5, \texttt{reg\_lambda}=2.0). Table~\ref{tab:overfit_initial} reveals severe overfitting, particularly in XGBoost.

\begin{table}[t]
\caption{Overfitting Diagnosis --- Initial Model (5-Fold CV)}
\label{tab:overfit_initial}
\centering
\small
\begin{tabular}{@{}lcccc@{}}
\toprule
\textbf{Model} & \textbf{Train ROC} & \textbf{Val ROC} & \textbf{Train PR} & \textbf{Val PR} \\
\midrule
XGBoost & \textbf{1.0000} & 0.9455 & \textbf{1.0000} & 0.7051 \\
LightGBM & 0.9893 & 0.9476 & 0.8809 & 0.6872 \\
CatBoost & 0.9853 & 0.9516 & 0.8799 & 0.7220 \\
\bottomrule
\end{tabular}
\vspace{-1mm}
\end{table}

\textbf{Key finding}: XGBoost achieves \textit{perfect} training scores (ROC-AUC = 1.0, PR-AUC = 1.0) but only 0.7051 validation PR-AUC---a gap of 0.295. This indicates the model memorizes the training set rather than learning generalizable patterns. The train-validation gap is the true measure of overfitting, not the validation score alone.

\subsection{Overfitting Diagnosis: Improved Model}

After applying aggressive regularization (Table~\ref{tab:params}), feature selection, and repeated CV, Table~\ref{tab:overfit_improved} shows dramatically reduced overfitting.

\begin{table}[t]
\caption{Overfitting Diagnosis --- Improved Model (3$\times$5-Fold CV)}
\label{tab:overfit_improved}
\centering
\small
\begin{tabular}{@{}lcccc@{}}
\toprule
\textbf{Model} & \textbf{Train ROC} & \textbf{Val ROC} & \textbf{Train PR} & \textbf{Val PR} \\
\midrule
XGBoost & 0.9947 & 0.9468 & 0.9410 & 0.7032 \\
LightGBM & 0.9649 & 0.9453 & 0.7050 & 0.6515 \\
CatBoost & 0.9681 & 0.9503 & 0.7785 & 0.7213 \\
\bottomrule
\end{tabular}
\vspace{-1mm}
\end{table}

Table~\ref{tab:gap_reduction} quantifies the gap reduction.

\begin{table}[t]
\caption{Train--Validation Gap Reduction}
\label{tab:gap_reduction}
\centering
\small
\begin{tabular}{@{}lcccc@{}}
\toprule
& \multicolumn{2}{c}{\textbf{ROC-AUC Gap}} & \multicolumn{2}{c}{\textbf{PR-AUC Gap}} \\
\textbf{Model} & Initial & Improved & Initial & Improved \\
\midrule
XGBoost & 0.055 & \textbf{0.048} & 0.295 & \textbf{0.238} \\
LightGBM & 0.042 & \textbf{0.020} & 0.194 & \textbf{0.054} \\
CatBoost & 0.034 & \textbf{0.018} & 0.158 & \textbf{0.057} \\
\bottomrule
\end{tabular}
\vspace{-1mm}
\end{table}

\textbf{LightGBM} achieves the largest improvement: PR-AUC gap drops from 0.194 to 0.054 (\textbf{72\% reduction}). CatBoost's gap drops 64\%. Both models are now considered healthy (gap $< 0.10$). XGBoost remains the most overfit (gap = 0.238) despite improvement, reflecting its aggressive fitting strategy.

\subsection{Model Comparison: Initial vs Improved}

Table~\ref{tab:comparison} compares the initial and improved ensemble on aggregated out-of-fold predictions.

\begin{table}[t]
\caption{Ensemble Performance: Initial vs Improved (OOF)}
\label{tab:comparison}
\centering
\small
\begin{tabular}{@{}lccccc@{}}
\toprule
\textbf{Model} & \textbf{ROC} & \textbf{PR} & \textbf{F1} & \textbf{Prec} & \textbf{Rec} \\
\midrule
Initial Ensemble & \textbf{0.9533} & 0.7010 & \textbf{0.6808} & \textbf{0.6384} & 0.7292 \\
Improved Ensemble & 0.9527 & \textbf{0.7042} & 0.6772 & 0.6075 & \textbf{0.7648} \\
\midrule
$\Delta$ & $-$0.0006 & +0.0032 & $-$0.0036 & $-$0.0309 & \textbf{+0.0356} \\
\bottomrule
\end{tabular}
\vspace{-1mm}
\end{table}

\textbf{Key observations}:
\begin{itemize}[leftmargin=*]
    \item \textbf{PR-AUC improved} from 0.7010 to 0.7042---the primary metric for imbalanced fraud detection.
    \item \textbf{Recall increased} from 72.9\% to 76.5\%, catching 18 additional fraudulent providers (387 vs 369 out of 506).
    \item \textbf{Precision decreased} from 63.8\% to 60.8\%---an expected tradeoff when optimizing for recall in fraud investigation pipelines.
    \item \textbf{ROC-AUC is essentially unchanged} ($-$0.0006), confirming the initial validation scores were trustworthy.
\end{itemize}

The improved model is preferable for deployment: it catches more fraud (higher recall) with more trustworthy estimates (smaller overfit gap), while the modest precision reduction is acceptable in investigation-oriented workflows where human auditors review flagged cases.

\subsection{Confusion Matrix Analysis}

\begin{table}[h]
\caption{Confusion Matrices (Out-of-Fold, $n$ = 5,410)}
\label{tab:confusion}
\centering
\small
\begin{tabular}{@{}l|cc|cc@{}}
\toprule
& \multicolumn{2}{c|}{\textbf{Initial ($\tau$=0.575)}} & \multicolumn{2}{c}{\textbf{Improved ($\tau$=0.528)}} \\
& Pred Neg & Pred Pos & Pred Neg & Pred Pos \\
\midrule
Actual Neg & 4,695 & 209 & 4,654 & 250 \\
Actual Pos & 137 & 369 & 119 & 387 \\
\bottomrule
\end{tabular}
\vspace{-1mm}
\end{table}

The improved model catches 18 more fraudulent providers (387 vs 369) at the cost of 41 additional false positives. In operational terms: 18 fewer fraud cases are missed, while 41 more legitimate providers are flagged for audit review---a favorable tradeoff when audit capacity exists.

\subsection{Threshold Analysis}

\begin{table}[h]
\caption{Threshold Analysis at Target Recall Levels (Improved)}
\label{tab:thresholds}
\centering
\small
\begin{tabular}{@{}cccc@{}}
\toprule
\textbf{Target Recall} & $\tau$ & \textbf{Precision} & \textbf{F1} \\
\midrule
70\% & 0.528 & 0.608 & 0.677 \\
80\% & 0.494 & 0.579 & 0.673 \\
90\% & 0.449 & 0.466 & 0.614 \\
\bottomrule
\end{tabular}
\vspace{-1mm}
\end{table}

At the operationally critical 80\% recall target, the improved model achieves 57.9\% precision with F1 = 0.673, meaning approximately 4 out of 5 fraudulent providers are detected while just over half the flagged cases are true positives.

\subsection{Adversarial Validation}

An XGBoost classifier trained to distinguish train from test feature vectors achieves ROC-AUC = 0.488 (3-fold CV), which is indistinguishable from random chance. This confirms that the train/test split does not introduce distribution shift, and performance estimates on training data should transfer to the test set.

\subsection{Calibration Assessment}

\begin{table}[h]
\caption{Probability Calibration Quality}
\label{tab:calibration}
\centering
\small
\begin{tabular}{@{}lcc@{}}
\toprule
\textbf{Model} & \textbf{Brier Score} $\downarrow$ & \textbf{ECE} $\downarrow$ \\
\midrule
XGBoost (improved) & \textbf{0.0590} & \textbf{0.0628} \\
Ensemble (calibrated) & 0.0624 & 0.1493 \\
\bottomrule
\end{tabular}
\vspace{-1mm}
\end{table}

The XGBoost base model shows better calibration (Brier = 0.059, ECE = 0.063) than the isotonic-calibrated ensemble (ECE = 0.149). This suggests that isotonic calibration on the small dataset introduces some distortion. In production, we recommend using the uncalibrated XGBoost probabilities for risk scoring and the ensemble's binary predictions for flag/no-flag decisions.

\subsection{Feature Importance Analysis}

\begin{table}[h]
\caption{Top 15 Features by XGBoost Gain (Improved Model)}
\label{tab:importance}
\centering
\small
\begin{tabular}{@{}clc@{}}
\toprule
\textbf{Rank} & \textbf{Feature} & \textbf{Gain} \\
\midrule
1 & MaxDiagCodes & 0.1901 \\
2 & InpatientClaims & 0.1022 \\
3 & MaxProcCodes & 0.0890 \\
4 & MaxSameDayClaimsPerBene & 0.0638 \\
5 & MaxClaimDuration & 0.0607 \\
6 & MaxReimbursed & 0.0424 \\
7 & MeanSameDayClaimsPerBene & 0.0423 \\
8 & TotalClaims & 0.0306 \\
9 & MeanClaimDuration & 0.0299 \\
10 & \textit{graph\_diag\_concentration} & 0.0246 \\
11 & \textit{graph\_proc\_concentration} & 0.0230 \\
12 & ClaimsPerBeneficiary & 0.0162 \\
13 & FracDead & 0.0123 \\
14 & MeanProcCodes & 0.0119 \\
15 & ClaimDateSpanDays & 0.0110 \\
\bottomrule
\end{tabular}
\vspace{-1mm}
\end{table}

\textbf{Key observations}:
\begin{itemize}[leftmargin=*]
    \item \textbf{MaxDiagCodes} remains dominant (19.0\%) but with reduced concentration compared to the initial model (27.3\%), indicating the improved regularization distributes importance more evenly across features---a sign of reduced overfitting.
    \item \textbf{InpatientClaims} rises to rank 2 (10.2\%) in the improved model, suggesting that the regularized model relies more on robust volume-based features.
    \item \textbf{Graph features} (ranks 10--11) maintain their position, confirming that network-derived concentration indices provide genuine signal rather than overfitting-driven importance.
    \item \textbf{FracDead} (rank 13) captures a known fraud pattern: billing for deceased beneficiaries.
\end{itemize}

\subsection{Individual Base Learner Comparison (Improved)}

\begin{table}[h]
\caption{Improved Base Learner Performance (3$\times$5-Fold CV)}
\label{tab:base_learners}
\centering
\small
\begin{tabular}{@{}lcccc@{}}
\toprule
\textbf{Model} & \textbf{Val ROC} & \textbf{Val PR} & \textbf{ROC Gap} & \textbf{PR Gap} \\
\midrule
XGBoost & 0.9468 & 0.7032 & 0.048 & 0.238 \\
LightGBM & 0.9453 & 0.6515 & \textbf{0.020} & \textbf{0.054} \\
CatBoost & \textbf{0.9503} & \textbf{0.7213} & 0.018 & 0.057 \\
\midrule
\textbf{Ensemble} & 0.9527 & \textbf{0.7042} & --- & --- \\
\bottomrule
\end{tabular}
\vspace{-1mm}
\end{table}

CatBoost achieves the best individual performance (ROC-AUC = 0.9503, PR-AUC = 0.7213) with the lowest overfitting gap. LightGBM has the smallest gaps but lower absolute performance. The ensemble combines their complementary strengths.

\subsection{Validation Stability}

The repeated CV design reveals fold-level variance:

\begin{table}[h]
\caption{Fold-Level Stability (Std across 15 evaluations)}
\label{tab:stability}
\centering
\small
\begin{tabular}{@{}lcc@{}}
\toprule
\textbf{Model} & \textbf{Val ROC Std} & \textbf{Val PR Std} \\
\midrule
XGBoost & 0.0097 & 0.0359 \\
LightGBM & 0.0114 & 0.0484 \\
CatBoost & 0.0103 & 0.0377 \\
\bottomrule
\end{tabular}
\vspace{-1mm}
\end{table}

ROC-AUC is highly stable (std $\approx$ 0.01) across all models. PR-AUC shows higher variance (std $\approx$ 0.04), expected given the small positive class ($\approx$101 fraud cases per fold).

% ============================================================
% V. DISCUSSION
% ============================================================
\section{Discussion}
\label{sec:discussion}

\subsection{The Importance of Train-Score Reporting}

Our most practically significant finding is that \textbf{reporting only validation metrics hides critical overfitting information}. The initial model's validation ROC-AUC of 0.9452 and the improved model's 0.9527 are nearly identical, yet the underlying model health differs dramatically: the initial XGBoost has a 0.295 PR-AUC gap while the improved LightGBM has only 0.054. We advocate for mandatory reporting of train--validation gaps in healthcare fraud detection publications.

\subsection{Regularization as the Primary Lever}

On a 5,410-row dataset, \textbf{regularization matters more than model architecture}. The single most impactful change was reducing \texttt{max\_depth} from 6 to 4, which prevents trees from memorizing provider-specific patterns. The combination of shallower trees, higher minimum leaf sizes (30--40), stronger L1/L2 penalties, and aggressive column subsampling (0.6) reduced the LightGBM PR gap by 72\%.

\subsection{Value of Graph Features}

Graph features contribute consistent improvement across both the initial and improved models. The HHI concentration indices (ranks 10--11) survive feature importance pruning and maintain their rank despite aggressive regularization, confirming they capture genuine fraud signal. The shared-beneficiary features (\texttt{graph\_total\_shared\_benes}) appear at rank 20, encoding provider collusion patterns inaccessible to independent tabular features.

\subsection{Limitations}

\begin{enumerate}[leftmargin=*]
    \item \textbf{Label quality}: Provider fraud labels in this dataset are binary without indication of fraud type or severity. In production, labels derived from OIG LEIE exclusion records would be more reliable.
    \item \textbf{Temporal validation}: The train/test split is not explicitly temporal. Production models should be trained on historical data and evaluated on future periods.
    \item \textbf{Dataset scale}: With 5,410 training providers, the dataset is small by modern ML standards, inherently limiting model complexity and validation reliability.
    \item \textbf{Residual XGBoost overfitting}: Despite aggressive regularization, XGBoost retains a 0.238 PR-AUC gap. This is inherent to the small dataset and XGBoost's fitting strategy; further reduction would require more training data.
    \item \textbf{Calibration distortion}: Isotonic calibration on the ensemble increased ECE from 0.063 to 0.149, likely due to the small dataset not supporting reliable non-parametric calibration.
\end{enumerate}

\subsection{Future Work}

\begin{itemize}[leftmargin=*]
    \item \textbf{Multi-year CMS data}: Extending to CMS PUF data (2013--2019) with LEIE-derived labels \cite{herland2018} would provide 10--50$\times$ more providers, enabling deeper models and richer graph features.
    \item \textbf{Temporal validation}: True out-of-time evaluation (train on years 1--3, test on year 4) would assess real-world deployment performance.
    \item \textbf{Graph neural networks}: On larger provider networks, GNN-based approaches could capture higher-order neighborhood structures beyond our hand-crafted graph features.
    \item \textbf{SHAP-based explainability}: Per-provider SHAP explanations would support auditor decision-making in operational deployment.
    \item \textbf{Platt calibration}: Sigmoid (Platt) calibration may outperform isotonic calibration on small datasets.
\end{itemize}

\subsection{Practical Deployment Considerations}

For production deployment, we recommend:
\begin{itemize}[leftmargin=*]
    \item \textbf{Threshold}: Use the 80\% recall threshold ($\tau$=0.494, precision=57.9\%) as the initial operating point.
    \item \textbf{Risk scoring}: Use uncalibrated XGBoost probabilities (Brier=0.059) for continuous risk scores; use ensemble binary predictions for flag/no-flag decisions.
    \item \textbf{Monitoring}: Track train--validation gaps across retraining cycles; rising gaps indicate model degradation.
    \item \textbf{Retraining}: Retrain quarterly on rolling 3-year windows.
\end{itemize}

% ============================================================
% VI. CONCLUSION
% ============================================================
\section{Conclusion}
\label{sec:conclusion}

We presented a production-grade, graph-enhanced ensemble learning pipeline for healthcare provider fraud detection with explicit overfitting diagnosis and mitigation. By combining 49 carefully selected features (from 58 tabular + 12 graph) in a calibrated stacking ensemble of XGBoost, LightGBM, and CatBoost with aggressive regularization, we achieved a PR-AUC of 0.7042 and recall of 76.5\% on 5,410 Medicare providers with a 9.35\% fraud rate.

The key contributions are: (1) demonstration that train--validation gap analysis is essential for healthcare fraud models, revealing hidden overfitting invisible in validation-only reporting; (2) a systematic regularization strategy that reduces the PR-AUC gap by up to 72\% while improving recall; (3) confirmation that Neo4j-derived graph features (concentration indices, shared-entity overlap) provide genuine, regularization-robust fraud signal; and (4) a comprehensive diagnostic framework including adversarial validation, calibration assessment, and fold-level stability analysis.

Future work should focus on scaling to multi-year CMS PUF datasets with LEIE-derived labels, enabling temporal validation and richer graph features that would further improve both absolute performance and generalization confidence.

% ============================================================
% REFERENCES
% ============================================================
\begin{thebibliography}{00}

\bibitem{nhcaa2020}
National Health Care Anti-Fraud Association, ``The challenge of health care fraud,'' NHCAA White Paper, 2020.

\bibitem{bauder2017}
R. A. Bauder and T. M. Khoshgoftaar, ``Medicare fraud detection using machine learning methods,'' in \textit{Proc. 16th IEEE Int. Conf. Machine Learning and Applications (ICMLA)}, 2017, pp. 858--865.

\bibitem{herland2018}
M. Herland, T. M. Khoshgoftaar, and R. A. Bauder, ``Big Data fraud detection using multiple medicare data sources,'' \textit{Journal of Big Data}, vol. 5, no. 1, pp. 1--21, 2018.

\bibitem{johnson2023}
J. M. Johnson and T. M. Khoshgoftaar, ``Medicare fraud detection using neural networks,'' \textit{SN Computer Science}, vol. 4, no. 2, pp. 1--13, 2023.

\bibitem{bauder2023}
R. A. Bauder and T. M. Khoshgoftaar, ``Multirelational data mining for healthcare fraud detection,'' \textit{SN Computer Science}, vol. 4, 2023.

\bibitem{pourhabibi2020}
T. Pourhabibi, K. L. Ong, B. H. Kam, and Y. L. Boo, ``Fraud detection: A systematic literature review of graph-based anomaly detection approaches,'' \textit{Decision Support Systems}, vol. 133, 2020.

\bibitem{liu2021}
C. Liu, Y. Weng, and M. Luo, ``Healthcare fraud detection via bipartite graph neural network,'' in \textit{Proc. IEEE Int. Conf. Bioinformatics and Biomedicine (BIBM)}, 2021.

\bibitem{vanvlasselaer2017}
V. Van Vlasselaer, T. Eliassi-Rad, L. Akoglu, M. Snoeck, and B. Baesens, ``GOTCHA! Network-based fraud detection for social security fraud,'' \textit{Management Science}, vol. 63, no. 9, pp. 3090--3110, 2017.

\bibitem{neo4j2023}
Neo4j, Inc., ``Neo4j Graph Database Platform,'' 2023. [Online]. Available: https://neo4j.com

\bibitem{wolpert1992}
D. H. Wolpert, ``Stacked generalization,'' \textit{Neural Networks}, vol. 5, no. 2, pp. 241--259, 1992.

\bibitem{chen2016}
T. Chen and C. Guestrin, ``XGBoost: A scalable tree boosting system,'' in \textit{Proc. 22nd ACM SIGKDD Int. Conf. Knowledge Discovery and Data Mining}, 2016, pp. 785--794.

\bibitem{ke2017}
G. Ke et al., ``LightGBM: A highly efficient gradient boosting decision tree,'' in \textit{Advances in Neural Information Processing Systems (NeurIPS)}, vol. 30, 2017.

\bibitem{prokhorenkova2018}
L. Prokhorenkova, G. Gusev, A. Vorobev, A. V. Dorogush, and A. Gulin, ``CatBoost: Unbiased boosting with categorical features,'' in \textit{Advances in Neural Information Processing Systems (NeurIPS)}, vol. 31, 2018.

\bibitem{niculescu2005}
A. Niculescu-Mizil and R. Caruana, ``Predicting good probabilities with supervised learning,'' in \textit{Proc. 22nd Int. Conf. Machine Learning (ICML)}, 2005, pp. 625--632.

\bibitem{vandenheuvel2022}
M. van den Heuvel, ``On the limitations of SMOTE for imbalanced learning on aggregated features,'' \textit{Machine Learning}, vol. 111, pp. 3519--3543, 2022.

\end{thebibliography}

\end{document}
